
\section{截面的\BookRecordIndex{T0337}伪循环}\label{sec:8-5}

本节从截面\BookRecordIndex{T0246}模空间 $\mathcal M(\widetilde A;J,H)$ 出发，在 $\widetilde M^k$ 中构造
\BookRecordIndex{T0337}伪循环。所需的许多概念与第~\ref{sec:6-7} 节十分接近。不过，我们现在的视角略有不同：
第~\ref{sec:6-7} 节的设置强调 $M$ 中的曲线 $u$ 而非其图像 $\widetilde u$，这里则完全用
$\widetilde u$ 来讨论。

固定\BookRecordIndex{T0306}竖直殆复结构 $J\in\mathcal J^{\mathrm{Vert}}$、满足
\[
  \pi_*\widetilde A=[\Sigma]
\]
的同调类 $\widetilde A\in H_2(\widetilde M;\mathbb Z)$，以及正则\BookRecordIndex{T0157}哈密顿扰动
$H\in\mathcal H_{\mathrm{reg}}(\widetilde A,J)$。这些数据确定表示类 $\widetilde A$ 的
$\widetilde J_H$-全纯截面所成的光滑\BookRecordIndex{T0246}模空间
$\mathcal M(\widetilde A;J,H)$。记
\[
  \mathcal M_{\Sigma,k}(\widetilde A;J,H)
  :=\bigl\{(\widetilde u,z_1,\ldots,z_k)\mid
  z_i\ne z_j\ \forall i\ne j\bigr\}
  \subset\mathcal M(\widetilde A;J,H)\times\Sigma^k.
\]
它的维数为
\[
  \dim\mathcal M_{\Sigma,k}(\widetilde A;J,H)
  =\mu(\widetilde A,g,k)
  :=n(2-2g)+2c_1^{\mathrm{Vert}}(\widetilde A)+2k,
\]
其中 $c_1^{\mathrm{Vert}}\in H^2(\widetilde M;\mathbb Z)$ 是\BookRecordIndex{T0308}竖直切丛的第一 Chern 类，
$g$ 是 $\Sigma$ 的亏格（见定理~\ref{thm:8-3-1}）。我们的目标是找出条件，使\BookRecordIndex{T0283}求值映射
\[
  \widetilde{\operatorname{ev}}
  =\widetilde{\operatorname{ev}}_{J,H}:
  \mathcal M_{\Sigma,k}(\widetilde A;J,H)\longrightarrow\widetilde M^k,
\]
\[
  \widetilde{\operatorname{ev}}(\widetilde u,z_1,\ldots,z_k)
  :=\bigl(\widetilde u(z_1),\ldots,\widetilde u(z_k)\bigr),
\]
成为\BookRecordIndex{T0337}伪循环。

我们需要对纤维施加一个比\BookRecordIndex{T0058}半正性略强的条件。这个条件必须保证：在正则情形，纤维中的
所有 $J$-全纯球面都有非负 Chern 数。\footnote{第~\ref{sec:6-7} 节不需要这个更强的
假设，因为在那里可以假定族 $\{J_z\}$ 几乎常值；比较
$\mathcal J_+(S^2;M,\omega;\kappa)$ 的定义。}由于已经假设基底 $\Sigma$ 连通，所有纤维
都辛同胚于同一个流形 $(M,\omega)$；对 $(M,\omega)$ 的相关条件是：它\BookRecordIndex{T0056}半单调，或者其第一
Chern 类在球面上为零，或者其最小 Chern 数 $N\ge n-1$。这个条件也可以写成
\begin{equation}\label{eq:8-5-1}
  \omega(A)>0,\qquad c_1(A)\ge2-n
  \quad\Longrightarrow\quad c_1(A)\ge0
\end{equation}
对每个 $A\in\pi_2(M)$ 都成立；这里
$c_1:=c_1(TM,J)\in H^2(M;\mathbb Z)$。它蕴含：定义~\ref{def:6-4-1} 中的\BookRecordIndex{T0058}半正性条件对
由 $(\widetilde M,\widetilde\omega)$ 各纤维内球面表示的类成立。若 $\Sigma$ 的亏格为正，
则只要 $\kappa:\Sigma\to\mathbb R$ 满足式~\eqref{eq:8-1-5}，
$(\widetilde M,\widetilde\omega_\kappa)$ 就是\BookRecordIndex{T0058}半正\BookRecordIndex{T0358}辛流形。不过，当 $\Sigma=S^2$ 时，
条件~\eqref{eq:8-5-1} 比 $\widetilde M$ 的\BookRecordIndex{T0058}半正性弱。第~\ref{sec:9-6} 节的主要应用所涉及的是
二维球面上的\BookRecordIndex{T0156}哈密顿纤维化，其中纤维单调，但总空间不\BookRecordIndex{T0058}半正。

\begin{theorem}\label{thm:8-5-1}
设 $(\widetilde M,\pi,\widetilde\omega)$ 是亏格为 $g$ 的紧黎曼曲面
$(\Sigma,j_\Sigma,\mathrm{dvol}_\Sigma)$ 上的\BookRecordIndex{T0200}局部哈密顿纤维化，其紧纤维维数为 $2n$，
并满足式~\eqref{eq:8-5-1}。固定整数 $k\ge0$ 和满足
$\pi_*\widetilde A=[\Sigma]$ 的同调类
$\widetilde A\in H_2(\widetilde M;\mathbb Z)$。则存在\BookRecordIndex{T0300}剩余集
\[
  \mathcal J\mathcal H_{\mathrm{reg}}
  (\widetilde M,\pi,\widetilde\omega)
  \subset\mathcal J^{\mathrm{Vert}}\times\mathcal H,
\]
具有以下性质。
\begin{enumerate}[label=\textup{(\roman*)},leftmargin=3em]
\item 对每个
$(J,H)\in\mathcal J\mathcal H_{\mathrm{reg}}
(\widetilde M,\pi,\widetilde\omega)$，\BookRecordIndex{T0283}求值映射
\[
  \widetilde{\operatorname{ev}}_{J,H}:
  \mathcal M_{\Sigma,k}(\widetilde A;J,H)\longrightarrow\widetilde M^k
\]
是维数为
$n(2-2g)+2c_1^{\mathrm{Vert}}(\widetilde A)+2k$ 的\BookRecordIndex{T0337}伪循环。
\item 若
$(J_0,H_0),(J_1,H_1)\in
\mathcal J\mathcal H_{\mathrm{reg}}(\widetilde M,\pi,\widetilde\omega)$，则
$\widetilde{\operatorname{ev}}_{J_0,H_0}$ 与
$\widetilde{\operatorname{ev}}_{J_1,H_1}$ 配边。
\end{enumerate}
\end{theorem}

集合
$\mathcal J\mathcal H_{\mathrm{reg}}(\widetilde M,\pi,\widetilde\omega)$ 的元素满足
第~\ref{sec:6-7} 节定义~\ref{def:6-7-10} 中条件的类似版本。精确表述见下文
定义~\ref{def:8-5-2}。

\paragraph{\BookRecordIndex{T0340}稳定映射。}
适当的\BookRecordIndex{T0340}稳定映射概念与第~\ref{sec:6-7} 节基本相同。固定\BookRecordIndex{T0306}竖直殆复结构
$J\in\mathcal J^{\mathrm{Vert}}$ 和\BookRecordIndex{T0157}哈密顿扰动 $H\in\mathcal H$。再固定一棵带 $k$ 个
标号、具有\BookRecordIndex{T0321}特殊顶点 $0\in T$ 的树 $T=(T,E,\Lambda)$，以及一族同调类
$\{\widetilde A_\alpha\}_{\alpha\in T}$，其中根类 $\widetilde A_0$ 是\BookRecordIndex{T0189}截面类，
$\alpha\ne0$ 时 $\widetilde A_\alpha$ 是\BookRecordIndex{T0309}竖直球面类；\sourcefootnote{原书此处及本节后文
多次把整族 $\{\widetilde A_\alpha\}$ 统称为“球面同调类”。但式~(8.5.2) 规定
$\pi_*\widetilde A_0=[\Sigma]$，且根分支
$\widetilde u_0:\Sigma\to\widetilde M$ 是截面；只有 $\alpha\ne0$ 的类满足
$\pi_*\widetilde A_\alpha=0$ 并由纤维中的球面表示。故后文均按根类为\BookRecordIndex{T0189}截面类、其余为
\BookRecordIndex{T0309}竖直球面类表述。}假设它们满足稳定性条件~\eqref{eq:6-1-1} 以及
\begin{equation}\label{eq:8-5-2}
  \pi_*\widetilde A_0=[\Sigma],
  \qquad
  \pi_*\widetilde A_\alpha=0
  \quad(\alpha\in T\setminus\{0\}).
\end{equation}
与这些数据相联系的\BookRecordIndex{T0340}稳定映射\BookRecordIndex{T0246}模空间为
\[
  \mathcal M_{\Sigma,T}(\{\widetilde A_\alpha\};J,H)
  :=\widetilde{\mathcal M}_{\Sigma,T}
  (\{\widetilde A_\alpha\};J,H)/G_T,
\]
定义如下。$\widetilde{\mathcal M}_{\Sigma,T}
(\{\widetilde A_\alpha\};J,H)$ 的一个元素是\BookRecordIndex{T0340}稳定映射
\[
  (\widetilde{\mathbf u},\mathbf z)
  =\bigl(\{\widetilde u_\alpha\}_{\alpha\in T},
    \{z_{\alpha\beta}\}_{\alpha E\beta},
    \{z_i\}_{1\le i\le k}\bigr),
\]
其中每个 $\widetilde u_\alpha:\Sigma_\alpha\to\widetilde M$ 都是表示类
$\widetilde A_\alpha$ 的 $\widetilde J_H$-全纯曲线；对每个 $\alpha\in T$，满足
$\alpha E\beta$ 的点 $z_{\alpha\beta}\in\Sigma_\alpha$，连同满足
$\alpha_i=\alpha$ 的点 $z_i\in\Sigma_\alpha$，两两不同；并且当 $\alpha E\beta$ 时
\[
  \widetilde u_\alpha(z_{\alpha\beta})
  =\widetilde u_\beta(z_{\beta\alpha}).
\]
此外，还要求
\[
  \Sigma_0=\Sigma,
  \qquad \pi\circ\widetilde u_0=\operatorname{id}_\Sigma,
  \qquad \Sigma_\alpha=\mathbb{CP}^1\quad(\alpha\ne0).
\]
这就是把\BookRecordIndex{T0340}稳定映射概念自然推广到伪全纯截面。\BookRecordIndex{T0413}重参数化群 $G_T$ 的元素是元组
$(f,\{\phi_\alpha\}_{\alpha\in T})$，其中 $f:T\to T$ 是满足 $f(0)=0$ 的树自同构，
$\phi_0=\operatorname{id}$，且当 $\alpha\ne0$ 时
$\phi_\alpha\in G=\operatorname{PSL}(2,\mathbb C)$。它像定义~\ref{def:5-1-4} 中那样作用
在\BookRecordIndex{T0340}稳定映射\BookRecordIndex{T0246}模空间上。当 $\Sigma=\mathbb{CP}^1$ 时，这只是第~5 章定义的一个特例。

定理~\ref{thm:8-5-1} 只涉及所有\BookRecordIndex{T0066}标记点均可自由变化的截面\BookRecordIndex{T0246}模空间的性质。所得\BookRecordIndex{T0337}伪循环定义
\BookRecordIndex{T0066}标记点不受限制的\BookRecordIndex{T0143}格罗莫夫--威滕不变量。命题 8.6.5 将说明：选取 $\widetilde M$ 中
Poincaré 对偶可由某个纤维的子流形表示的上同调类，便能恢复\BookRecordIndex{T0150}固定标记点的不变量。不过，
为了计算这些不变量并证明相关\BookRecordIndex{T0256}粘合定理，有时让某些\BookRecordIndex{T0066}标记点固定更为方便。这促使我们考察
部分\BookRecordIndex{T0066}标记点固定的\BookRecordIndex{T0340}稳定映射\BookRecordIndex{T0246}模空间。（后续应用中只考虑两个极端情形：或者没有点固定，或者
所有点都固定。这里为完整起见纳入一般情形。）

固定子集 $I\subset\{1,\ldots,k\}$，以及由 $I$ 标号的 $\Sigma$ 上两两不同的点所成元组
$\mathbf w=\{w_i\}_{i\in I}\in\Sigma^I$。它确定子集
\[
  \mathcal M_{\Sigma,T}(\{\widetilde A_\alpha\};\mathbf w,J,H)
  :=\widetilde{\mathcal M}_{\Sigma,T}
    (\{\widetilde A_\alpha\};\mathbf w,J,H)/G_T,
\]
由所有还满足下列条件的\BookRecordIndex{T0340}稳定映射 $(\widetilde{\mathbf u},\mathbf z)$ 组成：
\begin{equation}\label{eq:8-5-3}
  i\in I
  \quad\Longrightarrow\quad
  \pi\bigl(\widetilde u_{\alpha_i}(z_i)\bigr)=w_i.
\end{equation}
相应的\BookRecordIndex{T0184}简单稳定映射\BookRecordIndex{T0246}模空间记为
\[
  \mathcal M^*_{\Sigma,T}(\{\widetilde A_\alpha\};\mathbf w,J,H)
  :=\widetilde{\mathcal M}^*_{\Sigma,T}
    (\{\widetilde A_\alpha\};\mathbf w,J,H)/G_T;
\]
其定义与第~\ref{sec:6-1} 节完全相同。特别地，
\begin{equation}\label{eq:8-5-4}
  \mathcal M_{\Sigma,k}(\widetilde A;\mathbf w,J,H)
  :=\mathcal M(\widetilde A;J,H)\times\Sigma^k(\mathbf w),
\end{equation}
其中 $\Sigma^k(\mathbf w)\subset\Sigma^k$ 是由 $\Sigma$ 上两两不同点构成的所有
$k$-元组 $\mathbf z=(z_1,\ldots,z_k)$ 之集，并满足 $i\in I$ 时 $z_i=w_i$。

与第~\ref{sec:6-2} 节一样，\BookRecordIndex{T0184}简单稳定映射之集可描述为\BookRecordIndex{T0283}求值映射的原像。有一个由表示各类
$\widetilde A_\alpha$ 的简单元组 $\{\widetilde u_\alpha\}$ 构成的\BookRecordIndex{T0246}模空间
$\mathcal M^*(\{\widetilde A_\alpha\};J,H)$。若对每个
$\alpha\in T\setminus\{0\}$ 都有
$J\in\mathcal J_{\mathrm{reg}}^{\mathrm{Vert}}(\widetilde A_\alpha)$，并且
$H\in\mathcal H_{\mathrm{reg}}(\widetilde A_0,J)$，则它是定向有限维流形（见
式~\eqref{eq:8-3-2}）。再以 $Z(T)$ 表示所有元组
\[
  \mathbf z=\bigl(\{z_{\alpha\beta}\}_{\alpha E\beta},
    \{z_i\}_{1\le i\le k}\bigr)
\]
之集，其中对每个 $\alpha\in T$，满足 $\alpha E\beta$ 的点
$z_{\alpha\beta}\in\Sigma_\alpha$，连同满足 $\alpha_i=\alpha$ 的点
$z_i\in\Sigma_\alpha$，两两不同。于是有\BookRecordIndex{T0283}求值映射
\begin{equation}\label{eq:8-5-5}
  \widetilde{\operatorname{ev}}^E\times\pi^I:
  \mathcal M^*(\{\widetilde A_\alpha\};J,H)\times Z(T)
  \longrightarrow\widetilde M^E\times\Sigma^I,
\end{equation}
定义为
\[
  \widetilde{\operatorname{ev}}^E(\widetilde{\mathbf u},\mathbf z)
  :=\{\widetilde u_\alpha(z_{\alpha\beta})\}_{\alpha E\beta},
  \qquad
  \pi^I(\widetilde{\mathbf u},\mathbf z)
  :=\{\pi(\widetilde u_{\alpha_i}(z_i))\}_{i\in I}.
\]
满足式~\eqref{eq:8-5-3} 的\BookRecordIndex{T0184}简单稳定映射\BookRecordIndex{T0246}模空间，是
$\widetilde{\operatorname{ev}}^E\times\pi^I$ 下子流形
$\widetilde\Delta^E\times\{\mathbf w\}$ 的原像。这引出下面的定义。

\begin{definition}\label{def:8-5-2}
设
$\mathbf w\in\Sigma^I$。\sourcefootnote{原书在定义 8.5.2 与定理 8.5.3 中均把此处的
$\Sigma^I$ 印作 $\Sigma_I$；但前文只定义了由 $I$ 标号点组成的乘积 $\Sigma^I$，而
式~\eqref{eq:8-5-5} 及下文普适投影的靶也都是 $\Sigma^I$，故两处统一补正为上标。}
若 $(J,H)\in\mathcal J^{\mathrm{Vert}}\times\mathcal H$ 满足下列条件，就称它
\emph{关于 $\mathbf w$ 正则}。
\begin{description}[leftmargin=3.8em,style=nextline]
\item[\textup{(H)}] 对每个满足 $\pi_*\widetilde A=[\Sigma]$ 的
$\widetilde A\in H_2(\widetilde M;\mathbb Z)$，都有
$H\in\mathcal H_{\mathrm{reg}}(\widetilde A,J)$。
\item[\textup{(V)}] 对每个满足 $\pi_*\widetilde A=0$ 的球面同调类
$\widetilde A\in H_2(\widetilde M;\mathbb Z)$，都有
$J\in\mathcal J_{\mathrm{reg}}^{\mathrm{Vert}}(\widetilde A)$。
\item[\textup{(EW)}] 对每一棵带 $k$ 个标号、具有\BookRecordIndex{T0321}特殊顶点 $0$ 的树
$T=(T,E,\Lambda)$，以及每一族满足稳定性条件~\eqref{eq:6-1-1} 和
式~\eqref{eq:8-5-2} 的同调类（根为\BookRecordIndex{T0189}截面类，其余为\BookRecordIndex{T0309}竖直球面类），\BookRecordIndex{T0283}求值映射
\eqref{eq:8-5-5} 都与 $\widetilde\Delta^E\times\{\mathbf w\}$ 横截。
\end{description}
关于 $\mathbf w$ 正则的二元组 $(J,H)$ 之集一般记为
$\mathcal J\mathcal H_{\mathrm{reg}}
(\widetilde M,\pi,\widetilde\omega;\mathbf w)$；当 $I=\varnothing$ 时记为
$\mathcal J\mathcal H_{\mathrm{reg}}
(\widetilde M,\pi,\widetilde\omega)$。
\end{definition}

\paragraph{主要结果的证明。}

\begin{theorem}\label{thm:8-5-3}
对每个 $\mathbf w\in\Sigma^I$，集合
$\mathcal J\mathcal H_{\mathrm{reg}}
(\widetilde M,\pi,\widetilde\omega;\mathbf w)$ 都是
$\mathcal J^{\mathrm{Vert}}\times\mathcal H$ 中的\BookRecordIndex{T0300}剩余集。
\end{theorem}

\begin{proof}[证明概述]
给定带 $k$ 个标号的\BookRecordIndex{T0297}森林 $T=(T,E,\Lambda)$ 和一族满足式~\eqref{eq:8-5-2} 的同调类
$\{\widetilde A_\alpha\}_{\alpha\in T}$，用
$\mathcal J\mathcal H^\ell_{\mathrm{reg}}
(T,\{\widetilde A_\alpha\};\mathbf w)$ 表示所有 $C^\ell$ 类二元组
\[
  (J,H)\in\mathcal J^{\mathrm{Vert},\ell}\times\mathcal H^\ell
\]
中对 $T$ 与 $\{\widetilde A_\alpha\}$ 满足定义~\ref{def:8-5-2} 要求者之集。\BookRecordIndex{T0213}可分四步
证明：对所有 $T$ 与 $\{\widetilde A_\alpha\}$，集合
$\mathcal J\mathcal H^\ell_{\mathrm{reg}}
(T,\{\widetilde A_\alpha\};\mathbf w)$ 在
$\mathcal J^{\mathrm{Vert},\ell}\times\mathcal H^\ell$ 中稠密。

第一步是观察到\BookRecordIndex{T0268}普适模空间
\[
  \mathcal M^*(\{\widetilde A_\alpha\};
    \mathcal J^{\mathrm{Vert},\ell},\mathcal H^\ell)
\]
是\BookRecordIndex{T0053}巴拿赫流形。这来自命题~\ref{prop:3-2-1}、命题~\ref{prop:6-2-7} 和
定理~\ref{thm:8-3-1} 证明中论证的组合。

第二步是证明：若\BookRecordIndex{T0269}普适求值映射
\[
  \widetilde{\operatorname{ev}}^{E,\mathcal J,\mathcal H}:
  \mathcal M^*(\{\widetilde A_\alpha\};
    \mathcal J^{\mathrm{Vert},\ell},\mathcal H^\ell)
  \times Z(T)\longrightarrow\widetilde M^E
\]
与 $\widetilde\Delta^E$ 横截，那么每个 $\mathbf w\in\Sigma^I$ 都是普适投影
\[
  \pi^{I,\mathcal J,\mathcal H}:
  \mathcal M^*_{\Sigma,T}
  (\{\widetilde A_\alpha\};
    \mathcal J^{\mathrm{Vert},\ell},\mathcal H^\ell)
  :=\bigl(\widetilde{\operatorname{ev}}^{E,\mathcal J,\mathcal H}\bigr)^{-1}
    (\widetilde\Delta^E)
  \longrightarrow\Sigma^I
\]
的\BookRecordIndex{T0409}正则值。这与定理~\ref{thm:6-7-11} 证明的步骤 2 相同。

第三步是证明\BookRecordIndex{T0283}求值映射
$\widetilde{\operatorname{ev}}^{E,\mathcal J,\mathcal H}$ 确实与
$\widetilde\Delta^E$ 横截。对\BookRecordIndex{T0297}森林的边数归纳即可；归纳步骤把第二步与
定理~\ref{thm:6-3-1} 证明中的论证组合起来，正如定理~\ref{thm:6-7-11} 证明的步骤 3。

第四步考察投影
\[
  \pi^\ell:
  \mathcal M^*_{\Sigma,T}
  (\{\widetilde A_\alpha\};\mathbf w,
    \mathcal J^{\mathrm{Vert},\ell},\mathcal H^\ell)
  :=\bigl(\pi^{I,\mathcal J,\mathcal H}\bigr)^{-1}(\mathbf w)
  \longrightarrow\mathcal J^{\mathrm{Vert},\ell}\times\mathcal H^\ell.
\]
这是弗雷德霍姆映射；二元组
$(J,H)\in\mathcal J^{\mathrm{Vert},\ell}\times\mathcal H^\ell$ 是 $\pi^\ell$ 的\BookRecordIndex{T0409}正则值，
当且仅当
$(J,H)\in\mathcal J\mathcal H^\ell_{\mathrm{reg}}
(T,\{\widetilde A_\alpha\};\mathbf w)$。因此 萨德--斯梅尔定理断言，后一集合在
$\mathcal J^{\mathrm{Vert},\ell}\times\mathcal H^\ell$ 中稠密。

像定理~\ref{thm:3-1-6} 的证明一样，用 Taubes 的论证把 $C^\infty$ 情形化归到
$C^\ell$ 情形。最后，对所有 $T$ 与 $\{\widetilde A_\alpha\}$ 取可数交，便知
$\mathcal J\mathcal H_{\mathrm{reg}}
(\widetilde M,\pi,\widetilde\omega;\mathbf w)$ 是
$\mathcal J^{\mathrm{Vert}}\times\mathcal H$ 中的\BookRecordIndex{T0300}剩余集。这就证明了
定理~\ref{thm:8-5-3}。
\end{proof}

\begin{proof}[定理~\ref{thm:8-5-1} 的证明]
证明与定理~\ref{thm:6-7-1} 类似，这里只概述要点。取
$(J,H)\in\mathcal J\mathcal H_{\mathrm{reg}}
(\widetilde M,\pi,\widetilde\omega)$。对任意带 $k$ 个标号、具有\BookRecordIndex{T0321}特殊顶点 $0\in T$ 的树
$T$，以及任意一族满足稳定性条件~\eqref{eq:6-1-1} 与式~\eqref{eq:8-5-2} 的同调类
$\{\widetilde B_\alpha\}_{\alpha\in T}$（根为\BookRecordIndex{T0189}截面类，其余为\BookRecordIndex{T0309}竖直球面类），考察\BookRecordIndex{T0283}求值映射
\[
  \widetilde{\operatorname{ev}}_T:
  \mathcal M^*_{\Sigma,T}(\{\widetilde B_\alpha\})
  \longrightarrow\widetilde M^k,
\]
\[
  \widetilde{\operatorname{ev}}_T(\widetilde{\mathbf u},\mathbf z)
  :=\bigl(\widetilde u_{\alpha_1}(z_1),\ldots,
    \widetilde u_{\alpha_k}(z_k)\bigr).
\]
简单的维数计数给出
\begin{equation}\label{eq:8-5-6}
  \dim\mathcal M^*_{\Sigma,T}(\{\widetilde B_\alpha\})
  =\mu(\widetilde B,g,k)-2e(T),
  \qquad
  \widetilde B:=\sum_{\alpha\in T}\widetilde B_\alpha.
\end{equation}
此外，由式~\eqref{eq:8-5-1}，除非
\[
  c_1(\widetilde B_\alpha)\ge0
  \qquad(\alpha\in T\setminus\{0\}),
\]
否则这个\BookRecordIndex{T0246}模空间为空。\sourcefootnote{原书把此非负性写成对每个 $\alpha\in T$ 都成立；
但条件~(8.5.1) 只约束纤维中的\BookRecordIndex{T0285}球面类，而根类满足
$\pi_*\widetilde B_0=[\Sigma]$，是\BookRecordIndex{T0189}截面类，其竖直 Chern 数不必非负。故唯一可\BookRecordIndex{T0326}推出的范围是
$\alpha\ne0$。下文投影次数之和为一还蕴含根类重数为一，所以这一补正不影响随后所需的
维数比较。}

现在，由\BookRecordIndex{T0139}格罗莫夫紧性（当 $\Sigma=\mathbb{CP}^1$ 时是定理~\ref{thm:5-3-1}，高亏格
情形作显然推广）和命题~\ref{prop:6-1-2}，\BookRecordIndex{T0283}求值映射
$\widetilde{\operatorname{ev}}$ 的极限集可由上述\BookRecordIndex{T0283}求值映射
$\widetilde{\operatorname{ev}}_T$ 的像之并覆盖；这里遍历所有具有\BookRecordIndex{T0321}特殊顶点 $0$ 的带
$k$ 个标号树 $T$，以及所有满足稳定性条件~\eqref{eq:6-1-1}、式~\eqref{eq:8-5-2}、
$e(T)>0$ 和
\begin{equation}\label{eq:8-5-7}
  \widetilde A
  =\sum_{\alpha\in T}m_\alpha\widetilde B_\alpha
\end{equation}
的同调类族，其中 $\{m_\alpha\}_{\alpha\in T}$ 是正整数族。

关键观察是，由\BookRecordIndex{T0141}格罗莫夫收敛的重标度公理，极限\BookRecordIndex{T0340}稳定映射中出现的每个\BookRecordIndex{T0275}气泡都取值于
$\widetilde M$ 的单个纤维。也可以这样论证：对极限\BookRecordIndex{T0340}稳定映射中的每条
$\widetilde J_H$-全纯曲线
$\widetilde u_\alpha:\Sigma_\alpha\to\widetilde M$，投影
$\pi\circ\widetilde u_\alpha:\Sigma_\alpha\to\Sigma$ 是全纯的，故次数非负。由于
$\pi_*\widetilde A=[\Sigma]$，这些次数之和为一。因此，其中一个
$\pi\circ\widetilde u_\alpha$ 的次数为一，其余次数为零。次数为一的分支是单纯的，
故底\BookRecordIndex{T0072}层\BookRecordIndex{T0184}简单稳定映射也具有同样性质。现在，维数公式~\eqref{eq:8-5-6} 表明所有这些\BookRecordIndex{T0246}模空间
的维数至多为 $\mu(\widetilde A,g,k)-2$。这证明了 (i)。

现在令
$(J_i,H_i)\in\mathcal J\mathcal H_{\mathrm{reg}}
(\widetilde M,\pi,\widetilde\omega)$（$i=0,1$），并考察相应的\BookRecordIndex{T0283}求值映射
\[
  \widetilde{\operatorname{ev}}_0:
  \mathcal M_{\Sigma,k}(\widetilde A;J_0,H_0)\to\widetilde M^k,
  \qquad
  \widetilde{\operatorname{ev}}_1:
  \mathcal M_{\Sigma,k}(\widetilde A;J_1,H_1)\to\widetilde M^k.
\]
我们证明，在定义~\ref{def:6-5-1} 的意义下，
$\widetilde{\operatorname{ev}}_1$ 与
$\widetilde{\operatorname{ev}}_0$ 配边。这里只概述所需的同伦论证。选取同伦
\[
  [0,1]\longrightarrow\mathcal J^{\mathrm{Vert}}\times\mathcal H,
  \qquad \lambda\longmapsto(J_\lambda,H_\lambda),
\]
使它满足下列条件。
\begin{description}[leftmargin=3.8em,style=nextline]
\item[\textup{(H)}] 对每个满足 $\pi_*\widetilde A=[\Sigma]$ 的
$\widetilde A\in H_2(\widetilde M;\mathbb Z)$，都有
$\{H_\lambda\}\in
\mathcal H_{\mathrm{reg}}(\widetilde A,\{J_\lambda\};H_0,H_1)$。
\item[\textup{(V)}] 对每个满足 $\pi_*\widetilde A=0$ 的球面同调类
$\widetilde A\in H_2(\widetilde M;\mathbb Z)$，都有
$\{J_\lambda\}\in
\mathcal J_{\mathrm{reg}}^{\mathrm{Vert}}(\widetilde A;J_0,J_1)$。
\item[\textup{(EW)}] 对每一棵带 $k$ 个标号、具有\BookRecordIndex{T0321}特殊顶点 $0$ 的树
$T=(T,E,\Lambda)$，以及每一族满足稳定性条件~\eqref{eq:6-1-1} 和
式~\eqref{eq:8-5-2} 的同调类（根为\BookRecordIndex{T0189}截面类，其余为\BookRecordIndex{T0309}竖直球面类），\BookRecordIndex{T0283}求值映射
\[
  \widetilde{\operatorname{ev}}^E:
  \mathcal W^*(\{\widetilde A_\alpha\};\{J_\lambda,H_\lambda\})
  \times Z(T)\longrightarrow\widetilde M^E
\]
都与 $\widetilde\Delta^E$ 横截。
\end{description}
这种\BookRecordIndex{T0407}正则同伦的存在性与定理~\ref{thm:8-5-3} 的证明相同。选定\BookRecordIndex{T0407}正则同伦后，考察\BookRecordIndex{T0246}模空间
\[
  \mathcal W_{\Sigma,k}
  (\widetilde A;\{J_\lambda,H_\lambda\})
  :=\bigl\{(\lambda,\widetilde u,\mathbf z)\mid0\le\lambda\le1,\
  (\widetilde u,\mathbf z)
  \in\mathcal M_{\Sigma,k}(\widetilde A;J_\lambda,H_\lambda)\bigr\}.
\]
由定理~\ref{thm:8-3-3}，这是带边光滑定向流形，且
\[
  \partial\mathcal W_{\Sigma,k}
  (\widetilde A;\{J_\lambda,H_\lambda\})
  =\mathcal M_{\Sigma,k}(\widetilde A;J_1,H_1)
   \cup\bigl(-\mathcal M_{\Sigma,k}(\widetilde A;J_0,H_0)\bigr).
\]
它的维数为 $\mu(\widetilde A,g,k)+1$，并带有\BookRecordIndex{T0283}求值映射
\begin{equation}\label{eq:8-5-8}
  \widetilde{\operatorname{ev}}:
  \mathcal W_{\Sigma,k}
  (\widetilde A;\{J_\lambda,H_\lambda\})
  \longrightarrow\widetilde M^k,
\end{equation}
\[
  \widetilde{\operatorname{ev}}
  (\lambda,\widetilde u,z_1,\ldots,z_k)
  :=\bigl(\widetilde u(z_1),\ldots,\widetilde u(z_k)\bigr).
\]
再次由\BookRecordIndex{T0139}格罗莫夫紧性和命题~\ref{prop:6-1-2}，\BookRecordIndex{T0283}求值映射~\eqref{eq:8-5-8} 的极限集可由
\BookRecordIndex{T0283}求值映射
\[
  \widetilde{\operatorname{ev}}_T:
  \mathcal W^*_{\Sigma,T}
  (\{\widetilde B_\alpha\};\{J_\lambda,H_\lambda\})
  \longrightarrow\widetilde M^k
\]
的像之并覆盖；这里遍历所有带 $k$ 个标号的树 $T$ 和所有满足稳定性条件
~\eqref{eq:6-1-1}、式~\eqref{eq:8-5-2}、$e(T)>0$ 及式~\eqref{eq:8-5-7} 的同调类族
$\{\widetilde B_\alpha\}_{\alpha\in T}$。这个并集中的\BookRecordIndex{T0246}模空间维数为
\[
  \mu(\widetilde B,g,k)+1-2e(T)
  \le\mu(\widetilde A,g,k)-1.
\]
因此\BookRecordIndex{T0283}求值映射~\eqref{eq:8-5-8} 是\BookRecordIndex{T0337}伪循环，并给出从
$\widetilde{\operatorname{ev}}_0$ 到
$\widetilde{\operatorname{ev}}_1$ 的所需配边。这就完成了定理~\ref{thm:8-5-1} 的证明
概述。
\end{proof}

下面的习题在某些\BookRecordIndex{T0066}标记点固定时定义类似的\BookRecordIndex{T0337}伪循环。

\begin{exercise}\label{ex:8-5-4}
固定整数 $k\ge0$ 和子集 $I\subset\{1,\ldots,k\}$。$\Sigma$ 上两两不同点所成的任意
元组 $\mathbf w=\{w_i\}_{i\in I}$ 确定乘积 $\widetilde M^k$ 的子流形
\[
  \widetilde M^k(\mathbf w)
  :=\bigl\{\widetilde{\mathbf x}
  =(\widetilde x_1,\ldots,\widetilde x_k)\in\widetilde M^k\mid
  \pi(\widetilde x_i)=w_i\text{ 对 }i\in I\bigr\},
\]
以及式~\eqref{eq:8-5-4} 中的\BookRecordIndex{T0246}模空间
$\mathcal M_{\Sigma,k}(\widetilde A;\mathbf w,J,H)$。证明：若
\[
  (J,H)\in\mathcal J\mathcal H_{\mathrm{reg}}
  (\widetilde M,\pi,\widetilde\omega;\mathbf w),
\]
则\BookRecordIndex{T0283}求值映射
\begin{equation}\label{eq:8-5-9}
  \widetilde{\operatorname{ev}}_{\mathbf w}
  =\widetilde{\operatorname{ev}}_{\mathbf w,J,H}:
  \mathcal M_{\Sigma,k}(\widetilde A;\mathbf w,J,H)
  \longrightarrow\widetilde M^k(\mathbf w),
\end{equation}
\[
  \widetilde{\operatorname{ev}}_{\mathbf w}
  (\widetilde u,z_1,\ldots,z_k)
  :=\bigl(\widetilde u(z_1),\ldots,\widetilde u(z_k)\bigr)
\]
是维数为
\[
  n(2-2g)+2c_1^{\mathrm{Vert}}(\widetilde A)+2k-2|I|
\]
的\BookRecordIndex{T0337}伪循环。再设 $\widetilde{\operatorname{ev}}_I$ 是
$\widetilde{\operatorname{ev}}_{\mathbf w}$ 与包含映射
$\widetilde M^k(\mathbf w)\to\widetilde M^k$ 的复合，\BookRecordIndex{T0326}推出映射
\begin{equation}\label{eq:8-5-10}
  \widetilde{\operatorname{ev}}_I:
  \mathcal M_{\Sigma,k}(\widetilde A;\mathbf w,J,H)
  \longrightarrow\widetilde M^k
\end{equation}
也是同一维数的\BookRecordIndex{T0337}伪循环。
\end{exercise}
